\documentclass[11pt,a4paper]{article}
\usepackage[margin=1in]{geometry}
\usepackage[T1]{fontenc}
\usepackage[utf8]{inputenc}
\usepackage{graphicx}
\usepackage{hyperref}
\usepackage{enumitem}
\usepackage{float}
\usepackage{amsmath}

\title{Advanced Image Processing Methods\\Homework 2 Report}
\author{Peter Dobosi}
\date{\today}

\begin{document}
\maketitle

\section*{Goal}
The task is to locate ROI (Region of Interest) query images inside a noisy full-size motherboard photograph.
The pipeline uses convolutional preprocessing, local feature detection, and descriptor matching to find each ROI's position and mark it with a bounding rectangle.

\section*{Method Summary}
\begin{enumerate}[leftmargin=*]
    \item \textbf{Preprocessing (Subtask~1)}:
        Both the full scene image and every ROI image are converted to grayscale and denoised with a $5\times5$ Gaussian blur.
        An unsharp-mask sharpening pass is applied after denoising (strength $0.30$) to recover fine texture detail lost during blurring, which significantly improves descriptor distinctiveness (see \textit{Parameter Tuning} below).
        ROI loading is dynamic using \texttt{Path.glob}, so the pipeline is independent of the number of ROI files.
    \item \textbf{Keypoint \& Descriptor Extraction (Subtask~2)}:
        ORB (\texttt{cv2.ORB\_create}, \texttt{nfeatures=10000}) was chosen over SIFT for the following reasons:
        \begin{itemize}[leftmargin=1.5em]
            \item Binary BRIEF-based descriptors are matched with Hamming distance, which is significantly faster than SIFT's 128-float L2 matching on a $2962\times1935$ scene image.
            \item The ROIs are axis-aligned crops at the same scale as the scene, so SIFT's scale-space advantage is unnecessary.
            \item ORB is the detector demonstrated in the course lab slides (slide~04) and matches the recommended BFMatcher workflow.
        \end{itemize}
        Keypoints and descriptors are computed for the denoised scene image and for each denoised ROI image.
    \item \textbf{Feature Matching (Subtask~3)}:
        A \texttt{cv2.BFMatcher} with \texttt{NORM\_HAMMING} and \texttt{crossCheck=False} is used with \texttt{knnMatch(k=2)}.
        Weak matches are filtered using the Lowe ratio test (threshold $0.75$).
    \item \textbf{Localization \& Visualization (Subtask~4)}:
        The centre of each ROI is estimated from the matched scene keypoint coordinates using a robust outlier-rejection scheme based on the Median Absolute Deviation (MAD).
        The median of the matched coordinates serves as an outlier-resistant initial estimate; points beyond $2\times\text{MAD}$ from the median are discarded, and the final centre is the mean of the remaining inliers.
        This mitigates the effect of false-positive matches that would otherwise shift the na\"ive mean estimate (see \textit{Observations \& Robustness} below).
        Bounding boxes are drawn with \texttt{cv2.rectangle} using the ROI dimensions, clamped to the scene boundaries.
        Per-ROI debug figures (side-by-side keypoints and match lines via \texttt{cv2.drawMatches}) are saved alongside the final localization image.
\end{enumerate}

\section*{Parameter Tuning}
The initial configuration (no sharpening, \texttt{nfeatures=3000}) produced very few good matches for the two smaller ROIs.
I iteratively refined two parameters and measured good matches after Lowe ratio filtering ($0.75$):

\begin{center}
\begin{tabular}{l|c c c c}
    \textbf{Configuration} & \textbf{roi\_0} & \textbf{roi\_1} & \textbf{roi\_2} & \textbf{roi\_3} \\
    \hline
    No sharpen, 3k features & 1 & 6 & 137 & 58 \\
    Sharpen 0.30, 3k features & 2 & 2 & 142 & 121 \\
    \textbf{Sharpen 0.30, 10k features} & \textbf{7} & \textbf{11} & \textbf{203} & \textbf{184} \\
\end{tabular}
\end{center}

Sharpening boosted matches for the larger ROIs by enhancing fine texture, while increasing \texttt{nfeatures} to 10\,000 was essential for the smaller ROIs: the $2962\times1935$ scene image needs a dense keypoint grid so that even small query regions have nearby matches.

\section*{Observations \& Robustness}
The full scene image is a blue GigaByte motherboard PCB placed upside-down (rotated $180^{\circ}$ clockwise) on an annotated checkerboard mat with taped measurement marks, on a brown table.
The four ROIs capture mostly unpopulated RAM and PCI-Express slots and their retention clips --- regions that contain several visually similar alternatives across the board.

With the final \texttt{nfeatures=10000} configuration, the two larger ROIs (\texttt{roi\_2} and \texttt{roi\_3}) produce 203 and 184 good matches respectively, yielding highly accurate bounding boxes.
The two smaller ROIs (\texttt{roi\_0} with 7 matches and \texttt{roi\_1} with 11 matches) are roughly half the pixel area of the other two, which limits the number of detectable keypoints.
With so few matches, even a single false-positive ``long line'' (a match whose scene-side keypoint lies far from the true ROI location) noticeably shifts a na\"ive mean estimate.

I considered three mitigation strategies:
\begin{enumerate}[leftmargin=*]
    \item \textbf{Tighter Lowe ratio threshold} --- reducing from $0.75$ to, e.g., $0.60$ would keep only the most distinctive matches.
        However, with only 7--11 matches to begin with, stricter filtering risks dropping below the minimum needed for any reliable position estimate.
    \item \textbf{Outlier rejection (MAD-based)} --- computing the median of match coordinates first, then discarding points beyond $2\times$ the Median Absolute Deviation.
        The median is inherently robust to outliers, and after rejection the mean of the remaining inliers gives an accurate centre without losing ``borderline but valid'' matches.
    \item \textbf{Proximity-weighted averaging} --- weighting each match's contribution to the centre by the density of nearby matches (e.g., inverse distance to the median).
        This softly down-weights outliers instead of discarding them, but adds complexity and is sensitive to the choice of kernel bandwidth.
\end{enumerate}
I adopted strategy~2 (MAD outlier rejection) as it is simple, parameter-light (only the $2\times$ threshold), and directly removes the false positives visible in the debug figures while preserving all genuine matches.

\section*{Results}
\begin{figure}[H]
    \centering
    \includegraphics[width=\textwidth]{output/final_localization.png}
    \caption{Final localization result with ROI bounding rectangles drawn on the full scene image.}
    \label{fig:final}
\end{figure}

\section*{Conclusion \& Future Work}
The pipeline successfully localises all four ROIs in the noisy motherboard scene.
The larger ROIs (\texttt{roi\_2}, \texttt{roi\_3}) produce dense, accurate matches and tightly aligned bounding boxes.
The smaller ROIs (\texttt{roi\_0}, \texttt{roi\_1}) are more challenging: their content is dominated by homogeneous slot surfaces with little distinctive texture, so the detection algorithms lock onto whatever small differences they can find --- typically corners of nearby silkscreen letters that are close together.
Because these sparse keypoints cluster around incidental lettering rather than spreading across the full ROI, the resulting centre estimates are biased towards those letter corners.

Two directions could improve accuracy for such homogeneous ROIs in future work:
\begin{enumerate}[leftmargin=*]
    \item \textbf{Colour-guided refinement} --- after obtaining a rough position estimate from greyscale keypoint matching, the original BGR colour information could be used to refine the location.
        For example, a normalised cross-correlation or template-matching pass on the colour image within a neighbourhood of the initial estimate would exploit chromatic cues that the greyscale ORB pipeline discards.
    \item \textbf{Localised brute-force search} --- once the keypoint-based estimate places the ROI approximately, a bounded sliding-window or template match could be run within a limited region (e.g.\ $\pm$\,50\% of the ROI dimensions around the estimate).
        This trades additional computation for a more precise alignment, and the restricted search area keeps the cost manageable compared to a full-image brute-force scan.
\end{enumerate}

\section*{Reproducibility}
From the \texttt{homework\_2/} directory:
\begin{itemize}[leftmargin=*]
    \item Install dependencies: \texttt{pip install -e .}
    \item Run pipeline: \texttt{python src/mw79on\_submission\_hw2/main.py}
\end{itemize}
The script generates \texttt{output/final\_localization.png} and per-ROI match figures in \texttt{output/matches/}.
This report is provided as \texttt{REPORT.pdf}.

\end{document}
